\uuid{PvLl}
\titre{Second membre/condition initiale}
\auteur{hanani ; mzari}
\organisation{exo7}
\licence{CC BY-NC-SA 4.0}
\datecreate{2020-01-20}
\niveau{L1}
\module{Analyse}
\chapitre{Equation différentielle}
\sousChapitre{Résolution d'équation différentielle du premier ordre}
\difficulte{1}
\theme{Equations du premier ordre, second membre, condition initiale}
\exo7qcmid{lille-2.equations-differentielles.003}

\contenu{
\begin{question}
\qclassification{225.01}
\qtags{motcle=second membre/condition initiale}
On considère l'équation différentielle $(E)$ : $y'-y=\mathrm{e}^t$. Parmi les affirmations suivantes cocher celles qui sont vraies :
\begin{answers}
\good{La fonction $\displaystyle y_0=t\mathrm{e}^{t}$ est une solution de $(E)$.}
\bad{La fonction $\displaystyle y_1=\mathrm{e}^{-t}$ est une solution de $(E)$.}
\bad{La fonction $\displaystyle y=(1-t)\mathrm{e}^{t}$ est l'unique solution de $(E)$ telle que $y(1)=0$.}
\good{La fonction $\displaystyle y=(1+t)\mathrm{e}^{t}$ est l'unique solution de $(E)$ telle que $y(0)=1$.}
\end{answers}
\begin{explanations}
On vérifie que $y_0=t\mathrm{e}^{t}$ est une solution de $(E)$. On vérifie aussi que $\displaystyle y=(1+t)\mathrm{e}^{t}$ est une solution de $(E)$ et en plus $y(0)=1$. Cette dernière est donc l'unique solution de $(E)$ telle que $y(0)=1$.
\end{explanations}
\end{question}
}
